% GENERATED FILE. Edit canonical lessons, then run npm run generate:books.
% canonical-source-hash: fnv1a64:a605add97a33ef9d
% canonical-lessons: ZH-C423-lukuang, ZH-C423-dongnan, ZH-W423-nan-south, ZH-C423-xinan, ZH-C423-nanxia, ZH-C423-nanbei

\chapter{Going South}
\label{ch:zh-going-south}
\hlchaptermodality{423}

\begin{chapteropening}
\textbf{By the end of this chapter:} \emph{I can ask how the roads are, name the southeast and southwest, say we are heading south and compare.}

Complete the last lesson of chapter 423: I can ask how the roads are, name the southeast and southwest, say we are heading south and compare.
\end{chapteropening}

\section[lùkuàng]{lùkuàng — road conditions}

\label{lesson:ZH-C423-lukuang}



\begin{quote}
Before the new one: say the Mandarin for to head up north, then the Mandarin for high-speed.
\end{quote}

\subsection*{You'll want to know: lùkuàng}
\emph{lùkuàng} — \textquotedblleft{}road conditions\textquotedblright{}.

\begin{quote}
\emph{lùkuàng bù hǎo dehuà wǒmen zuò dìtiě qù} — If the roads are bad, we'll take the metro.
\end{quote}

Its characters come later; for now it is a word for the ear.

\begin{cousinweb}
\emph{lùkuàng} is built from \zh{路} 'a road' and \emph{kuàng} 'a state of things', both inherited from Middle Chinese, \emph{luH} and \emph{xjwangH} (Guangyun, in Baxter's transcription), which glosses the second 'to compare'; the sense 'state of things' is not given there.
\end{cousinweb}

\begin{grammarlens}[title={What you've built}]
One more word for getting around.
\end{grammarlens}

\subsection*{Your turn}
\begin{itemize}
\raggedright
  \item \emph{Say it:} \emph{lùkuàng}
  \item \emph{Say it:} \emph{lùkuàng}, once more
  \item \emph{Say it:} \emph{lùkuàng bù hǎo dehuà wǒmen zuò dìtiě qù} — If the roads are bad, we'll take the metro.
  \item \emph{Recall:} say the Mandarin for to head up north, then the Mandarin for high-speed, then say \emph{lùkuàng} again
\end{itemize}

\begin{tcolorbox}[breakable,colback=teal!4,colframe=teal!35!black,title={Before you move on}]
What does lùkuàng mean? (\textquotedblleft{}Road conditions\textquotedblright{}.) Say it once more.
\end{tcolorbox}
\section[\texorpdfstring{\zh{东南} (dōngnán)}{东南 (dōngnán)}]{\zh{东南} (dōngnán) — southeast}

\label{lesson:ZH-C423-dongnan}



\begin{quote}
Before the new one: say the Mandarin for high-speed, then the Mandarin for road conditions.
\end{quote}

\subsection*{You'll want to know: \zh{东南}}
\textbf{\zh{东南}} — \emph{dōngnán} — \textquotedblleft{}southeast\textquotedblright{}.

\begin{quote}
\emph{huǒchēzhàn zài jīchǎng de dōngnán fāngxiàng} — The railway station is southeast of the airport.
\end{quote}

The new character, in pieces: \zh{南} — nine strokes: \zh{十} on top; under it a frame open at the foot, a vertical down the left and a bar that turns down the right side and hooks; inside it a dot and a short stroke falling to the left, two bars, and a vertical down through them. The next lesson puts it in your hand.

\begin{cousinweb}
\textbf{\zh{东南}} (\emph{dōngnán}) is built from \zh{东} 'east' and \zh{南} 'south', both inherited from Middle Chinese, \emph{tuwng} and \emph{nom} (Guangyun, in Baxter's transcription), which glosses the second as the direction of fire.
\end{cousinweb}

\begin{grammarlens}[title={What you've built}]
One more word for getting around.
\end{grammarlens}

\subsection*{Your turn}
\begin{itemize}
\raggedright
  \item \emph{Say it:} \emph{dōngnán}
  \item \emph{Say it:} \emph{dōngnán}, once more
  \item \emph{Say it:} \emph{huǒchēzhàn zài jīchǎng de dōngnán fāngxiàng} — The railway station is southeast of the airport.
  \item \emph{Recall:} say the Mandarin for high-speed, then the Mandarin for road conditions, then say \emph{dōngnán} again
\end{itemize}

\begin{tcolorbox}[breakable,colback=teal!4,colframe=teal!35!black,title={Before you move on}]
What does \zh{东南} mean? (\textquotedblleft{}Southeast\textquotedblright{}.) Say it once more.
\end{tcolorbox}
\section[\texorpdfstring{\zh{南} (nán)}{南 (nán)}]{\zh{南} — \zh{十} over a frame}

\label{lesson:ZH-W423-nan-south}



\begin{quote}
The word first, then the character inside it.

\begin{itemize}
\raggedright
  \item \emph{Recall:} say \emph{dōngnán} — \textbf{R1}, one lesson back, before it decays
\end{itemize}

\begin{quote}
\zh{十} \zh{东}
\end{quote}

\emph{Read it:} each one aloud, with nothing said first — \emph{shí}, \emph{dōng}
\end{quote}

\subsection*{Script — \zh{十} over a frame}
\hlblockfigure{\detokenize{figures/ZH-W423-nan-south-filmstrip.pdf}}{How \zh{南} is written, stroke by stroke}

\begin{quote}
\zh{南}
\end{quote}

\textbf{Two pieces, top to bottom.}

\textbf{The top piece, \zh{十}, two strokes.} A short bar, then a vertical down through it.

\textbf{The frame, two strokes.} A vertical down the left. Then a bar across that turns down the right side and hooks to the left at the foot: the frame is open at the bottom.

\textbf{Inside, five strokes.} A dot and a short stroke falling to the left, side by side. Two bars, one under the other. Then a vertical from the upper bar down through the lower one.

2 + 2 + 5 = nine.

\subsection*{Your turn}
\begin{itemize}
\raggedright
  \item \emph{Write it:} \zh{南}, all nine, with the model in front of you
  \item \emph{Cover:} write it again from memory, and count the pieces, not the strokes
\end{itemize}

\begin{tcolorbox}[breakable,colback=teal!4,colframe=teal!35!black,title={Before you move on}]
How many strokes in \zh{南}? (\textbf{Nine}.) What sits on top? (\textbf{\zh{十}}.) Is the frame closed at the foot? (\textbf{No, it is open}.)

Source: \href{https://raw.githubusercontent.com/chanind/hanzi-writer-data/68d10a4b21150cae5e1ebbd223eed289cf32d90c/data/\%E5\%8D\%97.json}{Hanzi Writer Data \zh{南}, pinned stroke paths and medians}.
\end{tcolorbox}
\section[\texorpdfstring{\zh{西南} (xīnán)}{西南 (xīnán)}]{\zh{西南} (xīnán) — southwest; the Southwest (of China)}

\label{lesson:ZH-C423-xinan}



\begin{quote}
Before the new one: say the Mandarin for road conditions, then the Mandarin for the southeast.
\end{quote}

\subsection*{You'll want to know: \zh{西南}}
\textbf{\zh{西南}} — \emph{xīnán} — \textquotedblleft{}southwest; the Southwest (of China)\textquotedblright{}.

\begin{quote}
\emph{wǒmen wǎng xīnán kāi} — We're driving southwest.
\end{quote}

\begin{cousinweb}
\textbf{\zh{西南}} (\emph{xīnán}) is built from \zh{西} 'west' and \zh{南} 'south', both inherited from Middle Chinese, \emph{sej} and \emph{nom} (Guangyun, in Baxter's transcription).
\end{cousinweb}

\begin{grammarlens}[title={What you've built}]
One more word for getting around.
\end{grammarlens}

\subsection*{Your turn}
\begin{itemize}
\raggedright
  \item \emph{Say it:} \emph{xīnán}
  \item \emph{Say it:} \emph{xīnán}, once more
  \item \emph{Say it:} \emph{wǒmen wǎng xīnán kāi} — We're driving southwest.
  \item \emph{Recall:} say the Mandarin for road conditions, then the Mandarin for the southeast, then say \emph{xīnán} again
\end{itemize}

\begin{tcolorbox}[breakable,colback=teal!4,colframe=teal!35!black,title={Before you move on}]
What does \zh{西南} mean? (\textquotedblleft{}Southwest; the Southwest (of China)\textquotedblright{}.) Say it once more.
\end{tcolorbox}
\section[\texorpdfstring{\zh{南下} (nánxià)}{南下 (nánxià)}]{\zh{南下} (nánxià) — to go south; to head down south}

\label{lesson:ZH-C423-nanxia}



\begin{quote}
Before the new one: say the Mandarin for the southeast, then the Mandarin for the southwest.
\end{quote}

\subsection*{You'll want to know: \zh{南下}}
\textbf{\zh{南下}} — \emph{nánxià} — \textquotedblleft{}to go south; to head down south\textquotedblright{}.

\begin{quote}
\textbf{\zh{下个月我们开车南下}} \emph{xià ge yuè wǒmen kāichē nánxià} — Next month we're driving south.
\end{quote}

\begin{cousinweb}
\textbf{\zh{南下}} (\emph{nánxià}) is built from \zh{南} 'south' and \zh{下} 'to go down', both inherited from Middle Chinese, \emph{nom} and \emph{haeX} (Guangyun, in Baxter's transcription), which glosses the second 'low' and 'to come down'.
\end{cousinweb}

\begin{grammarlens}[title={What you've built}]
One more word for getting around.
\end{grammarlens}

\subsection*{Your turn}
\begin{itemize}
\raggedright
  \item \emph{Say it:} \emph{nánxià}
  \item \emph{Say it:} \emph{nánxià}, once more
  \item \emph{Say it:} \emph{xià ge yuè wǒmen kāichē nánxià} — Next month we're driving south.
  \item \emph{Recall:} say the Mandarin for the southeast, then the Mandarin for the southwest, then say \emph{nánxià} again
\end{itemize}

\begin{tcolorbox}[breakable,colback=teal!4,colframe=teal!35!black,title={Before you move on}]
What does \zh{南下} mean? (\textquotedblleft{}To go south; to head down south\textquotedblright{}.) Say it once more.
\end{tcolorbox}
\section[\texorpdfstring{\zh{南北} (nánběi)}{南北 (nánběi)}]{\zh{南北} (nánběi) — north and south}

\label{lesson:ZH-C423-nanbei}



\begin{quote}
Before the new one: say the Mandarin for the southwest, then the Mandarin for to head down south.
\end{quote}

\subsection*{You'll want to know: \zh{南北}}
\textbf{\zh{南北}} — \emph{nánběi} — \textquotedblleft{}north and south\textquotedblright{}.

\begin{quote}
\emph{Zhōngguó de nánběi hěn bù yíyàng} — The north and the south of China are very different.
\end{quote}

\begin{cousinweb}
\textbf{\zh{南北}} (\emph{nánběi}) is built from \zh{南} 'south' and \zh{北} 'north', both inherited from Middle Chinese, \emph{nom} and \emph{pok} (Guangyun, in Baxter's transcription).
\end{cousinweb}

\begin{grammarlens}[title={What you've built}]
One more word for getting around.
\end{grammarlens}

\subsection*{Your turn}
\begin{itemize}
\raggedright
  \item \emph{Say it:} \emph{nánběi}
  \item \emph{Say it:} \emph{nánběi}, once more
  \item \emph{Say it:} \emph{Zhōngguó de nánběi hěn bù yíyàng} — The north and the south of China are very different.
  \item \emph{Recall:} say the Mandarin for the southwest, then the Mandarin for to head down south, then say \emph{nánběi} again
\end{itemize}

\begin{tcolorbox}[breakable,colback=teal!4,colframe=teal!35!black,title={Before you move on}]
What does \zh{南北} mean? (\textquotedblleft{}North and south\textquotedblright{}.) Say it once more.
\end{tcolorbox}
